\subsection{Q1: is the crowd right?}\label{sec:res-q1}
% src: S-PAPER-3, S-PAPER-5, S-PAPER-10
The pre-registered rule gives this verdict: \VerdictQOne{}
The chi-square statistic is \QOneChiSq{} with \QOneDf{} degrees of freedom.
The p-value is $p=\QOneP$.
This p-value assumes independent markets, so it may be too small.
The effect is small and runs one way.
In \GapNegBinsMain{} of \GapBinsTotal{} bins the observed outcome-1 rate is below the mean price.
So outcome 1 is priced too high.
Take the \OutcomeOneYesN{} markets where outcome 1 is labeled Yes.
In \OutcomeOneYesBinsNegative{} of 9 bins (0 to 8), their rate is below the mean price.
The gaps are \OutcomeOneYesGapMin{} to \OutcomeOneYesGapMax{} percentage points.
This is descriptive, with no new test.
Over all \NMarketsKept{} markets the net gap is \GapNetDiffMain.
That is \GapNetYesShort{} fewer outcome-1 wins than the prices imply.
The largest per-bin gap is \LargestGapAbs{} (bin \LargestGapBin, counting from 0).
The largest chi-square contributor is bin \ChiTopBin{}, with \ChiTopShare{} of the chi-square.
We set no size line before the results, so small is our reading.
\Cref{fig:calib} shows the observed outcome-1 rate against the mean price in each bin.
\Cref{tab:bins} gives the numbers.

\begin{table}[H]
\centering
\caption{Calibration by price bin, with the gap as observed rate minus mean price in percent points. Why: the gap and its interval show where prices are off. Source: run \RunId.}
\label{tab:bins}
\small
\setlength{\tabcolsep}{4pt}
\begin{tabular}{lrrrcrr}
\toprule
Bin & $n$ & Mean price & Observed & \CiLevel{} CI & Gap & BH $p$ \\
\midrule
\binrow{0 to 10\%}{ZeroTen}\\ \midrule
\binrow{10 to 20\%}{TenTwenty}\\ \midrule
\binrow{20 to 30\%}{TwentyThirty}\\ \midrule
\binrow{30 to 40\%}{ThirtyForty}\\ \midrule
\binrow{40 to 50\%}{FortyFifty}\\ \midrule
\binrow{50 to 60\%}{FiftySixty}\\ \midrule
\binrow{60 to 70\%}{SixtySeventy}\\ \midrule
\binrow{70 to 80\%}{SeventyEighty}\\ \midrule
\binrow{80 to 90\%}{EightyNinety}\\ \midrule
\binrow{90 to 100\%}{NinetyHundred}\\
\bottomrule
\end{tabular}
\end{table}

The observed rate is below the mean price in every bin except the top one.
Without the minimum-trade filter the count is \GapNegBinsNoMin{} of \GapBinsTotal.
In equal-frequency bins it is \GapNegBinsEqFreq{} of \GapBinsTotal.
After BH correction, \QOneBinsBh{} of the \NBins{} bins are off ($p<\Alpha$).
After Bonferroni correction, \QOneBinsBonferroni{} bins are off.
The \CiLevel{} z interval leaves the mean price outside in \ZIntervalBinsOutsideN{} bins: \ZIntervalBinsOutsideList.
BH marks the bins \ZIntervalBinsBhOffList{} as off, which are the same bins.
For example, in the bin with mean price \BinSeventyEightyPrice{}, the observed outcome-1 rate is \BinSeventyEightyObs.
The gap is \BinSeventyEightyDiff{} (BH $p=\BinSeventyEightyBhP$).

The two cheapest bins are the longshot check.
In the cheapest bin the gap is \BinZeroTenDiff{} (BH $p=\BinZeroTenBhP$).
In the next bin the gap is \BinTenTwentyDiff{} (BH $p=\BinTenTwentyBhP$).
With the minimum-trade filter, cheap bets pay off less often than their price says.
The direction matches the horse-racing result \citep{snowberg2010favorite}.
Without the filter, the cheapest bin flips above the price.
Its gap is \NoMinBinZeroDiff{} (BH $p=\NoMinBinZeroPBh$).
The next bin has a gap of \NoMinBinOneDiff{} (BH $p=\NoMinBinOnePBh$).
Neither bin is off after BH.
So the longshot statement holds only with the filter.
Our gaps are per market and the dataset paper's are per trade.
So we compare the direction only, not the size.


\begin{figure}[H]
\centering
\figorbox{figures/fig4_calibration_gap}{\linewidth}
\caption{Observed outcome-1 rate against mean price per bin (top), and the per-bin gap in percentage points with \CiLevel{} z intervals (bottom). Why: the lower panel shows the gaps that the diagonal hides. Source: run \RunId{} and the local figure folder runs/local\_fig4gap\_v9.}
\label{fig:calib}
\end{figure}

\begin{figure}[H]
\centering
\figorbox{figures/fig5_ecdf_yes_no}{\linewidth}
\caption{ECDFs of \pOne{} for markets where outcome 1 won and markets where outcome 0 won. Why: the farther apart the ECDFs, the better the price separates the outcomes. Source: run \RunId.}
\label{fig:ecdf}
\end{figure}

Next come the score measures.
The mean squared residual of the price is \QOneBrier, the log loss is \QOneLogLoss, and the mean entropy is \QOneEntropy{} bits.
The outcome-1 rate among kept markets is \QOneYesRate, and the mean price is \PriceOneDayMean.
At the \FavoriteLine{} line, the classifier accuracy is \QOneAccuracy.
Its precision is \QOnePrecision{} and its recall is \QOneRecall.
The KS statistic between the prices of those two groups is $D=\QOneKsD$ (\Cref{fig:ecdf}).
Without the minimum-trade filter there are \QOneNoMinN{} markets.
The mean squared residual is then \QOneNoMinBrier{}, with $\chi^2=\QOneNoMinChiSq$ and $p=\QOneNoMinP$.
The chi-square test also rejects without the filter.
With \QOneEqFreqBins{} equal-frequency bins, BH marks \QOneEqFreqBinsBh{} bins as off.
The z interval leaves the price outside \ZIntervalEqFreqOutsideN{} bins, and BH marks \QOneEqFreqBinsBh.
They differ in the lowest bin, which has very few outcome-1 wins.
The z interval has no correction for the \QOneEqFreqBins{} bins.
Its spread comes from the observed rate, which is unreliable near zero.
